% =========================================================
% Hyperbolic Functions — Notes
% Chapter: elementary-functions
% Section: hyperbolic
% Volume III · Analysis
% =========================================================

\section{Hyperbolic Functions}
\label{sec:hyperbolic}

% ---------------------------------------------------------
% Toolkit
% ---------------------------------------------------------

\begin{tcolorbox}[title={Toolkit --- Hyperbolic Functions}]
\begin{itemize}
  \item $\cosh x = (e^x + e^{-x})/2$: even part of $e^x$;\;
    $\cosh x \geq 1$;\; $\cosh 0 = 1$.
  \item $\sinh x = (e^x - e^{-x})/2$: odd part of $e^x$;\;
    bijective $\mathbb{R} \to \mathbb{R}$.
  \item $e^x = \cosh x + \sinh x$ for all $x$.
  \item Hyperbolic Pythagorean: $\cosh^2 x - \sinh^2 x = 1$
    (note the minus sign, not plus).
  \item $(\sinh x)' = \cosh x$;\; $(\cosh x)' = \sinh x$
    (no sign change, unlike trigonometric).
  \item $\tanh x \in (-1,1)$ for all $x$;\;
    $\tanh x \to \pm 1$ as $x \to \pm\infty$.
  \item Inverse hyperbolic functions have explicit logarithmic formulas.
\end{itemize}
\end{tcolorbox}

% ---------------------------------------------------------
% Definition — Hyperbolic Functions
% ---------------------------------------------------------

\begin{tcolorbox}[title={Definition --- Hyperbolic Functions}]
\begin{definition}[Hyperbolic Functions]
\label{def:hyperbolic}
For all $x \in \mathbb{R}$:
\[
  \cosh x \;:=\; \frac{e^x + e^{-x}}{2},
  \qquad
  \sinh x \;:=\; \frac{e^x - e^{-x}}{2},
  \qquad
  \tanh x \;:=\; \frac{\sinh x}{\cosh x}.
\]
Additional: $\operatorname{sech} x := 1/\cosh x$,\;
$\operatorname{csch} x := 1/\sinh x$,\;
$\coth x := \cosh x / \sinh x$.
\end{definition}
\end{tcolorbox}

\begin{remark*}[Standard quantified statement]
\[
  \forall x \in \mathbb{R} :\;
  \cosh x = \frac{e^x + e^{-x}}{2} \geq 1,\;
  \sinh x = \frac{e^x - e^{-x}}{2},\;
  e^x = \cosh x + \sinh x.
\]
\end{remark*}

\begin{remark*}[Definition predicate reading]
\[
  \operatorname{HyperbolicCosh}(x,y)
  \;\colonequiv\;
  y = \frac{e^x + e^{-x}}{2}.
  \qquad
  \operatorname{HyperbolicSinh}(x,y)
  \;\colonequiv\;
  y = \frac{e^x - e^{-x}}{2}.
\]
\[
  \operatorname{HyperbolicTanh}(x,y)
  \;\colonequiv\;
  \cosh x \neq 0 \;\wedge\; y = \frac{\sinh x}{\cosh x}.
\]
\end{remark*}

\begin{remark*}[Negated quantified statement]
\[
  \neg\,\operatorname{HyperbolicTanh}(x,y)
  \;\iff\;
  y \neq \frac{\sinh x}{\cosh x}.
\]
Note $\cosh x > 0$ for all $x$, so $\operatorname{HyperbolicTanh}$
is defined on all of $\mathbb{R}$.
\end{remark*}

\begin{remark*}[Contrast with trigonometric functions]
\begin{center}
\begin{tabular}{@{}ll@{}}
\textbf{Trigonometric} & \textbf{Hyperbolic} \\
$\cos^2 x + \sin^2 x = 1$ & $\cosh^2 x - \sinh^2 x = 1$ \\
$(\cos x)' = -\sin x$ & $(\cosh x)' = \sinh x$ \\
$(\sin x)' = \cos x$ & $(\sinh x)' = \cosh x$ \\
$e^{ix} = \cos x + i\sin x$ & $e^x = \cosh x + \sinh x$ \\
Parametrises unit circle & Parametrises unit hyperbola \\
\end{tabular}
\end{center}
The minus sign in $\cosh^2 - \sinh^2 = 1$ versus $\cos^2 + \sin^2 =
1$ reflects the substitution $x \mapsto ix$ in Euler's formula.
\end{remark*}

% ---------------------------------------------------------
% Proposition — Hyperbolic Pythagorean Identity
% ---------------------------------------------------------

\begin{proposition}[Hyperbolic Pythagorean Identity]
\label{prop:hyperbolic-pythagorean}
For all $x \in \mathbb{R}$:
\[
  \cosh^2 x - \sinh^2 x = 1.
\]

\smallskip\noindent
\hyperref[prf:hyperbolic-pythagorean]{\textit{Go to proof.}}
\end{proposition}

\begin{remark*}[Standard quantified statement]
\[
  \forall x \in \mathbb{R} :\; \cosh^2 x - \sinh^2 x = 1.
\]
\end{remark*}

\begin{remark*}[Negated quantified statement]
\[
  \neg\,(\forall x:\cosh^2 x - \sinh^2 x=1)
  \;\iff\;
  \exists x\in\mathbb{R}:\cosh^2 x - \sinh^2 x\neq 1.
\]
The negation has no witness --- the identity holds universally.
\end{remark*}

\begin{remark*}[Definition predicate reading]
\[
  \operatorname{HyperbolicPythagorean}(x)
  \;\colonequiv\;
  \operatorname{HyperbolicCosh}(x,c)
  \;\wedge\; \operatorname{HyperbolicSinh}(x,s)
  \;\wedge\; c^2 - s^2 = 1.
\]
\end{remark*}

\begin{remark*}[Proof]
Direct computation:
\[
  \cosh^2 x - \sinh^2 x
  = \frac{(e^x+e^{-x})^2}{4} - \frac{(e^x-e^{-x})^2}{4}
  = \frac{4}{4} = 1.
\]
\end{remark*}

% ---------------------------------------------------------
% Proposition — Addition Formulas
% ---------------------------------------------------------

\begin{proposition}[Addition Formulas for Hyperbolic Functions]
\label{prop:hyperbolic-addition}
For all $x, y \in \mathbb{R}$:
\[
  \sinh(x+y) = \sinh x \cosh y + \cosh x \sinh y,
  \qquad
  \cosh(x+y) = \cosh x \cosh y + \sinh x \sinh y.
\]

\smallskip\noindent
\hyperref[prf:hyperbolic-addition]{\textit{Go to proof.}}
\end{proposition}

\begin{remark*}[Standard quantified statement]
\[
  \forall x,y\in\mathbb{R}:\;
  \sinh(x+y) = \sinh x\cosh y + \cosh x\sinh y
  \;\wedge\;
  \cosh(x+y) = \cosh x\cosh y + \sinh x\sinh y.
\]
\end{remark*}

\begin{remark*}[Negated quantified statement]
\[
  \neg\,(\text{addition formula for }\sinh)
  \;\iff\;
  \exists x,y\in\mathbb{R}:\;
  \sinh(x+y) \neq \sinh x\cosh y + \cosh x\sinh y.
\]
The negation has no witness.
\end{remark*}

\begin{remark*}[Definition predicate reading]
\[
  \operatorname{HyperbolicAddition}(x,y)
  \;\colonequiv\;
  \sinh(x+y) = \sinh x\cosh y + \cosh x\sinh y
  \;\wedge\;
  \cosh(x+y) = \cosh x\cosh y + \sinh x\sinh y.
\]
\end{remark*}

\begin{remark*}[Proof strategy]
Substitute the definitions and expand, or note the formulas follow
from $e^{x+y} = e^x e^y$ by taking even and odd parts.
\end{remark*}

\begin{remark*}[Comparison with trigonometric addition formulas]
The $\sinh$ formula matches $\sin(x+y)$ exactly. The $\cosh$ formula
differs from $\cos(x+y)$ only in the sign of the last term:
$+\sinh x\sinh y$ versus $-\sin x\sin y$. The sign difference is a
direct consequence of $\cosh^2 - \sinh^2 = 1$ versus
$\cos^2 + \sin^2 = 1$.
\end{remark*}

% ---------------------------------------------------------
% Proposition — Limits of tanh
% ---------------------------------------------------------

\begin{proposition}[Limits of $\tanh$]
\label{prop:tanh-limits}
\[
  \lim_{x \to +\infty} \tanh x = 1,
  \qquad
  \lim_{x \to -\infty} \tanh x = -1.
\]

\smallskip\noindent
\hyperref[prf:tanh-limits]{\textit{Go to proof.}}
\end{proposition}

\begin{remark*}[Standard quantified statement]
\[
  \operatorname{LimitAtInfinity}(\tanh x,\,+\infty,\,1)
  \;\wedge\;
  \operatorname{LimitAtInfinity}(\tanh x,\,-\infty,\,-1).
\]
\end{remark*}

\begin{remark*}[Negated quantified statement]
\[
  \neg\left(\lim_{x\to+\infty}\tanh x=1\right)
  \;\iff\;
  \exists\varepsilon_0>0,\;\forall M>0,\;\exists x>M:\;
  |\tanh x - 1|\geq\varepsilon_0.
\]
\end{remark*}

\begin{remark*}[Definition predicate reading]
\[
  \operatorname{LimitAtInfinity}(\tanh x,\,+\infty,\,1)
  \;\colonequiv\;
  \forall\varepsilon>0,\;\exists M\in\mathbb{R},\;\forall x:\;
  x>M \Rightarrow \operatorname{OutputTolerance}(\tanh x,1,\varepsilon).
\]
\end{remark*}

\begin{remark*}[Proof strategy]
Write $\tanh x = (1 - e^{-2x})/(1 + e^{-2x})$. As $x \to +\infty$,
$e^{-2x} \to 0$, giving $\tanh x \to 1$. The limit at $-\infty$
follows by oddness: $\tanh(-x) = -\tanh(x)$.
\end{remark*}

% ---------------------------------------------------------
% Definition — Inverse Hyperbolic Functions
% ---------------------------------------------------------

\begin{definition}[Inverse Hyperbolic Functions]
\label{def:inverse-hyperbolic}
Since $\sinh$ is strictly increasing and bijective $\mathbb{R} \to
\mathbb{R}$, its inverse $\operatorname{arcsinh}$ is defined on all
of $\mathbb{R}$. Since $\cosh$ is strictly increasing on
$[0,\infty)$ with range $[1,\infty)$, its inverse
$\operatorname{arccosh}$ is defined on $[1,\infty)$. Since $\tanh$
is strictly increasing with range $(-1,1)$, $\operatorname{arctanh}$
is defined on $(-1,1)$.

Explicit logarithmic formulas:
\begin{align*}
  \operatorname{arcsinh} x
  &= \ln\!\left(x + \sqrt{x^2+1}\right),
  \quad x \in \mathbb{R}. \\
  \operatorname{arccosh} x
  &= \ln\!\left(x + \sqrt{x^2-1}\right),
  \quad x \geq 1. \\
  \operatorname{arctanh} x
  &= \tfrac{1}{2}\ln\!\left(\frac{1+x}{1-x}\right),
  \quad |x| < 1.
\end{align*}
\end{definition}

\begin{remark*}[Standard quantified statement]
\[
  \forall x\in\mathbb{R}:\;
  \sinh(\operatorname{arcsinh} x)=x
  \;\wedge\;
  \operatorname{arcsinh}(\sinh x)=x.
\]
\[
  \forall x\geq 1:\;
  \cosh(\operatorname{arccosh} x)=x.
  \qquad
  \forall |x|<1:\;
  \tanh(\operatorname{arctanh} x)=x.
\]
\end{remark*}

\begin{remark*}[Definition predicate reading]
\[
  \operatorname{Arcsinh}(y,x)
  \;\colonequiv\;
  x\in\mathbb{R} \;\wedge\; \sinh x = y
  \;\colonequiv\;
  x = \ln\!\left(y+\sqrt{y^2+1}\right).
\]
\[
  \operatorname{Arctanh}(y,x)
  \;\colonequiv\;
  |y|<1 \;\wedge\; \tanh x = y
  \;\colonequiv\;
  x = \tfrac{1}{2}\ln\!\left(\frac{1+y}{1-y}\right).
\]
\end{remark*}

\begin{remark*}[Negated quantified statement]
\[
  \neg\,\operatorname{Arctanh}(y,x)
  \;\iff\;
  |y|\geq 1 \;\vee\; \tanh x \neq y.
\]
\end{remark*}

\begin{remark*}[Derivation of $\operatorname{arcsinh}$]
Set $y = \sinh x = (e^x - e^{-x})/2$ and let $u = e^x > 0$. Then
$y = (u - 1/u)/2$, so $u^2 - 2yu - 1 = 0$. The positive root is
$u = y + \sqrt{y^2+1}$, giving $x = \ln(y+\sqrt{y^2+1})$.
\end{remark*}

\begin{remark*}[Derivatives]
$(\operatorname{arcsinh} x)' = 1/\sqrt{x^2+1}$;\;
$(\operatorname{arccosh} x)' = 1/\sqrt{x^2-1}$ for $x > 1$;\;
$(\operatorname{arctanh} x)' = 1/(1-x^2)$ for $|x| < 1$.
These follow from the inverse function theorem applied to $\sinh$,
$\cosh$, and $\tanh$ respectively.
\end{remark*}
